\documentclass[conference]{IEEEtran}
\IEEEoverridecommandlockouts
% The preceding line is only needed to identify funding in the first footnote. If that is unneeded, please comment it out.
\usepackage{cite}
\usepackage{amsmath,amssymb,amsfonts}
\usepackage{algorithmic}
\usepackage{graphicx}
\usepackage{textcomp}
\usepackage{xcolor}
\def\BibTeX{{\rm B\kern-.05em{\sc i\kern-.025em b}\kern-.08em
    T\kern-.1667em\lower.7ex\hbox{E}\kern-.125emX}}

\newcommand{\TODO}[1]{\textbf{\textcolor{red}{[TODO: #1]}}}

\title{DQNGuard: Towards Open-World RF Preliminary-Action Detection}

\author{
\IEEEauthorblockN{
Adam Trott\IEEEauthorrefmark{*},
Cameron Popillo\IEEEauthorrefmark{\dag},
Nathaniel D. Bastian\IEEEauthorrefmark{\ddag},
Roulin Zhou\IEEEauthorrefmark{\dag},
and Gokhan Kul\IEEEauthorrefmark{*}
}
\IEEEauthorblockA{
\begin{tabular}{c}
\IEEEauthorrefmark{*}CIS and \IEEEauthorrefmark{\dag}Computer Engineering, University of Massachusetts Dartmouth, Dartmouth, MA, USA\\
\IEEEauthorrefmark{\ddag}Army Cyber Institute, United States Military Academy, West Point, NY, USA\\
Email: atrott@umassd.edu, cpopillo@umassd.edu, nathaniel.bastian@westpoint.edu, rzhou@umassd.edu, gkul@umassd.edu
\end{tabular}
}
}

\begin{document}

\maketitle

\begin{abstract}
RF sensing can provide early evidence for cyberspace windows of superiority, but closed-set RF classifiers are unsafe when adversaries emit signal behaviors outside the trained taxonomy. We study open-world RF preliminary-action detection: recognizing known RF behaviors while routing novel behaviors as unknown. We introduce DQNGuard, a budgeted open-set decision layer over a multi-domain preliminary-action backbone. The backbone maps OTA RF windows into IQ, FFT, DCT, and polar representations; DQNGuard then uses predicted-class calibration, variance/energy-style guard evidence, and known-only thresholding to decide whether each window should be retained as known preliminary-action evidence or routed to an unknown-behavior pool. We evaluate over-the-air WiFi, Bluetooth, and Zigbee captures spanning five behaviors: Scan, Burst, Sustain, Hop, and Replay. Under fixed Scan-surrogate known-budget calibration, DQNGuard improves unknown-class detection relative to VarMax and a DQN-IDS-style confidence head while maintaining low known-sample rejection. A Target--Surrogate Matrix further shows that surrogate-open calibration is highly target-dependent: withheld classes are not interchangeable proxies for future unknown behaviors. These results position DQNGuard as an RF sensing and triage layer for QR-CWoS pipelines, where known preliminary actions can support downstream ATT\&CK/EW reasoning and rejected unknowns can feed label-making and continual-learning workflows.

\end{abstract}

\begin{IEEEkeywords}
open-set recognition, RF signal classification, electronic warfare, preliminary actions, deep reinforcement learning, cyberspace windows of superiority
\end{IEEEkeywords}


\begin{figure*}[!t]
    \centering
    \includegraphics[width=\textwidth]{figures/hero_figure/hero_dqnguard_pipeline_s22_tikz.pdf}
    \caption{End-to-end DQNGuard evidence flow. OTA RF windows are mapped into IQ, FFT, DCT, and polar representations, fused, and processed by a preliminary-action CNN. DQNGuard uses predicted-class calibration, guard evidence, and a known-budget threshold to route each window to known preliminary-action evidence or to an unknown-behavior pool for downstream analysis.}
    \label{fig:dqnguard_pipeline}
\end{figure*}

\section{Introduction}
Modern military operations increasingly depend on decisions made across cyberspace and the electromagnetic spectrum (EMS). An RF behavior observed before, during, or around a cyber operation can provide early evidence that an adversary is preparing to scan, interfere, replay, sustain access, or shift emissions. In the Quantifiable and Resilient Cyberspace Windows of Superiority (QR-CWoS) setting, this early evidence matters because the value of a response depends on timing: a defender must recognize useful precursor behavior before the opportunity to act has closed. An RF sensing layer for QR-CWoS therefore cannot merely classify familiar emissions. It must also preserve uncertainty when an observed signal behavior falls outside the trained taxonomy.

This requirement conflicts with the closed-world assumption used by most RF classifiers. A standard deep classifier is trained to assign each input to one of the known classes, and its softmax output provides no native mechanism for saying that an over-the-air (OTA) RF window does not belong to any known behavior. In contested RF environments, this assumption is unsafe. Jammers, spoofers, replay emitters, cooperative radios, protocol variants, and hardware/channel effects can all produce behaviors that differ from the available training set. The dangerous failure mode is not simply that the classifier is wrong; it is that the classifier can be wrong with high confidence, routing a genuinely novel behavior into the nearest known class exactly when the downstream system should be alerted to ambiguity.

Open-set recognition (OSR) addresses this failure mode by allowing a model to retain known-class recognition while rejecting samples that do not conform to the learned decision structure. Prior confidence-based OSR work, including varMax~\cite{baye2024varmax,broggi2025varmax}, shows that classifier outputs such as top-score gaps, logit variance, and energy-style evidence can expose unfamiliar inputs without retraining the backbone. Trott et al.~\cite{trott2026model} extend this idea to RF modulation classification and show that RF novelty can be strongly dependent on the composition of the known and unknown classes. Tiwari et al.~\cite{tiwaridqn} further show that softmax-derived confidence features can support a DQN-style known/unknown decision head in network intrusion detection. Together, these results suggest that an OSR decision layer can serve as a useful front end for larger cyber/EMS reasoning systems. They also expose a key constraint for our setting: the decision layer must control how many known RF behaviors it rejects while still preserving candidate unknowns for later analysis.

We study this problem for RF preliminary actions (PAs). A preliminary action is an observable RF behavior that may precede, enable, or contextualize a later cyber, EMS, or cyber-physical effect. The evaluated PA classes are Scan, Burst, Sustain, Hop, and Replay. These labels are intentionally not final ATT\&CK or EW technique labels. Instead, they are precursor-level evidence: Scan-like behavior can support reconnaissance hypotheses, Replay-like behavior can support compromise or unauthorized-command hypotheses, and frequency-agile or sustained behavior can support other EW-relevant interpretations. Known PA detections can therefore feed downstream attack-chain or response-planning modules, while rejected unknowns form a behavior pool for LLM-assisted label-making and continual RF behavior discovery.

This paper focuses on the sensing and triage layer required before those downstream modules can act. We introduce DQNGuard, a budgeted OSR decision layer for OTA RF preliminary-action recognition. DQNGuard operates on a multi-domain PA backbone that maps RF windows into IQ, FFT, DCT, and polar representations and produces closed-set logits, softmax probabilities, intermediate features, and a predicted PA. The decision layer then applies three stages: predicted-class conditional calibration, variance/energy guard evidence, and known-only budgeted thresholding. The threshold is selected from known calibration samples so that unknown detection is constrained by an explicit known-rejection budget rather than by an unconstrained search for novelty.

We evaluate DQNGuard on OTA RF captures spanning WiFi, Bluetooth, and Zigbee emissions and five PA behaviors. The experiments compare DQNGuard against varMax and a DQN-IDS-style confidence head, then analyze how surrogate-open calibration transfers across target unknown classes. The goal is not to claim a complete QR-CWoS response system or a final semantic labeling system for all novel emissions. Rather, the goal is to determine whether an RF front end can separate known preliminary-action evidence from nonconforming unknown behavior under controlled calibration assumptions.

Our contributions are as follows:
\begin{enumerate}
    \item We formulate open-world RF preliminary-action detection as a sensing layer for QR-CWoS, where known PA detections provide precursor evidence and rejected unknowns provide a pool for downstream label-making.
    \item We introduce DQNGuard, a budgeted OSR decision layer that combines predicted-class calibration, variance/energy guard evidence, and known-only thresholding to control known-class rejection while identifying unknown RF behaviors.
    \item We evaluate DQNGuard on an OTA preliminary-action dataset spanning three protocol families and five behavior classes under fixed-surrogate, Target--Surrogate Matrix, and target-open calibration regimes.
    \item We show that surrogate-open calibration is target-dependent, motivating future work on principled surrogate selection, pooled calibration, and continual RF behavior discovery.
\end{enumerate}

The remainder of the paper is organized as follows. Section~\ref{sec:related} summarizes related work in OSR, RF signal classification, and DQN-based intrusion detection. Section~\ref{sec:methodology} defines the PA detection problem, baseline OSR heads, and DQNGuard. Section~\ref{sec:experiments} describes the OTA dataset, training splits, calibration regimes, and metrics. Section~\ref{sec:results} reports the open-set detection results. Section~\ref{sec:discussion} discusses operational implications for QR-CWoS, ATT\&CK/EW precursor evidence, and continual label-making, and Section~\ref{sec:conclusion} concludes.



\section{Related Works}
\label{sec:related}
Open-set recognition (OSR) addresses classification settings in which a model must preserve recognition of known classes while identifying samples that do not belong to any known class. This problem is especially important in cyber and electromagnetic sensing, where the operational value of a detector depends not only on assigning labels to familiar behaviors but also on refusing overconfident assignments for novel or adversarial behaviors. Confidence-based OSR methods provide a natural starting point because they operate on the decision structure of a closed-set classifier rather than requiring a separate generative model for every possible unknown.

The closest methodological predecessor to this work is varMax. Baye et al. introduced varMax as a confidence-based zero-day attack recognition method that uses maximum-score and variance-style evidence from classifier outputs to reject unknown samples~\cite{baye2024varmax}. Broggi et al. further develop the uncertainty and novelty-management perspective behind varMax, emphasizing the use of classifier confidence structure for open-set decisions~\cite{broggi2025varmax}. Our work inherits this view that novelty can be detected from the geometry of closed-set confidence outputs, but differs in two ways. First, DQNGuard conditions calibration on the predicted class rather than applying only a global rejection rule. Second, it uses a known-budget threshold so that unknown detection is constrained by an explicit limit on known-sample rejection.

A second line of related work uses reinforcement-learning-inspired decision layers for open-set intrusion detection. Tiwari et al. propose DQN-IDS, a deep reinforcement learning approach for open-set-enabled intrusion detection~\cite{tiwaridqn}. That work motivates treating OSR as a decision process over confidence states rather than as a pure feature-extraction problem. In this paper, we do not claim deployment-time continual reinforcement learning. Instead, we evaluate a DQN-IDS-style confidence head as a comparison point and use DQNGuard as a calibrated decision layer over fixed backbone outputs. This separation is useful experimentally because performance differences can be attributed to the OSR decision rule rather than to a different feature extractor or online learning policy.

Open-set RF and electromagnetic signal analysis has also explored multi-domain signal representations. Wei et al. exploit raw IQ, FFT, and DCT domains for detecting unclassified electromagnetic signals, using domain-specific autoencoders and feature fusion before open-set discrimination~\cite{wei2024multidomain}. Trott et al. similarly study RF modulation classifiers in the presence of novel samples and motivate multi-domain representations for RF OSR~\cite{trott2026model}. The present work adopts the same broad principle that RF novelty may appear differently across time-domain, frequency-domain, and transformed representations. However, rather than fusing autoencoder latents, our backbone processes IQ, FFT, DCT, and polar components through convolutional feature paths and then exposes logits, softmax probabilities, predicted labels, and intermediate features to the OSR decision layer.

This paper is therefore positioned at the intersection of three threads: confidence-based OSR, DQN-style open-set decision heads, and multi-domain RF feature extraction. The contribution is not a new RF waveform generator, a complete attack-chain planner, or an end-to-end autonomous EW response system. Instead, the paper asks whether a calibrated OSR decision layer can use multi-domain OTA RF preliminary-action evidence to separate known behaviors from novel behaviors under realistic calibration constraints. The resulting known detections can support precursor-level ATT\&CK/EW reasoning, while rejected unknowns can populate a downstream label-making and continual-learning loop.


\section{Methodology}
\label{sec:methodology}
\subsection{Threat Model and Operational Setting}
\label{subsec:threat-model}

We consider a contested Cyber-EMS environment in which a defender operates a sensor (a software-defined radio with antenna array) that observes a shared RF band. An adversary intermittently emits RF waveforms intended to enable a downstream electronic or cyber effect. For example, a burst pattern that probes a receiver's automatic gain control, a sustained emission used as a denial-of-service primitive, a hopping pattern used to evade fixed-channel monitoring, or a replay segment used to inject a previously captured cooperative signal. We refer to these emissions as \emph{preliminary actions} (PAs) following~\cite{trott2026model}, and we treat the \emph{full} attack which is the operational effect that the PA enables as occurring later in time.

Two assumptions are central. First, the adversary's PA repertoire is \emph{partially} known: the defender has labeled training examples for a finite set of PA classes, but the adversary may at any time deploy a PA whose waveform-level behavior falls outside that set. Second, the defender's response is \emph{costly} and \emph{time-bounded}: each candidate response (continued monitoring, additional sensing, deception, escalation) consumes resources and is only valuable if executed before the adversary's full tactic completes. The objective is therefore not classification accuracy in isolation, but the joint quantity \emph{probability of correctly characterizing the PA}~$\times$~\emph{duration of the resulting QR-CWoS}.

\subsection{Preliminary Action Taxonomy}
\label{subsec:pa-taxonomy}

Following the pilot taxonomy of~\cite{trott2026model}, we define four behaviorally distinct PA classes that serve as the trained label space:

\begin{itemize}
    \item \textbf{PA2: Burst.} The emitter transmits in many short on-intervals separated by quiet pauses; the defining feature is the temporal pattern of burst lengths and inter-burst gaps rather than the carrier itself.
    \item \textbf{PA3: Sustain.} The emitter transmits nearly continuously across the observation window with very few start/stop transitions, presenting as a steady occupied channel rather than discrete events.
    \item \textbf{PA4: Hop.} The emitter transmits while switching among a small set of frequencies, dwelling briefly on each and often revisiting previously used frequencies within the same window.
    \item \textbf{PA8: Replay.} The emitter records a short segment of a transmission and replays that exact segment repeatedly; replays are bit-for-bit copies of a waveform segment separated by short quiet gaps, which distinguishes them from newly generated bursts.
\end{itemize}

In addition, we hold out at least one PA class entirely from training to serve as the \emph{unseen-tactic} condition. The held-out PA is present in the over-the-air (OTA) test stream but never in training; the pipeline must label it \texttt{unknown} rather than collapse it into the nearest of \{PA2, PA3, PA4, PA8\}.

\subsection{Signal Acquisition and Preprocessing}
\label{subsec:preprocessing}

Each captured emission is recorded as a complex baseband I/Q stream and segmented into fixed-length analysis windows. For each window we compute (i)~a short-time Fourier spectrogram to expose the time-frequency structure that distinguishes burst from sustain and hop from replay, and (ii)~an event-detection trace that thresholds the per-bin energy to recover the on/off and dwell statistics on which PA2, PA3, and PA4 are partially defined. Windows are validated by the OTA validation pipeline of~\cite{trott2026model} to ensure TX/RX synchronization and to reject windows whose payload integrity has been compromised by channel effects; only validated windows enter the classifier. The output of preprocessing is a tensor representation $x \in \mathbb{R}^{C \times T \times F}$ together with a small vector of hand-engineered event-statistics features used as auxiliary input.

\subsection{Open-Set PA Classification with varMax}
\label{subsec:varmax}

The PA classifier is a deep neural network $f_\theta$ producing logits $z = f_\theta(x) \in \mathbb{R}^K$ over the $K=4$ trained PA classes. Rather than reading off $\arg\max \mathrm{softmax}(z)$, which is well known to be overconfident on out-of-distribution inputs, we apply the varMax open-set decision rule~\cite{baye2024varmax,broggi2025varmax}, which uses the \emph{variance} and the \emph{maximum} of the logit vector jointly as the basis for a confidence-aware decision. Concretely, at inference time we compute
\begin{equation}
\hat{y} =
\begin{cases}
\arg\max_k z_k & \text{if } g(z) \geq \tau \\
\texttt{unknown} & \text{otherwise}
\end{cases}
\label{eq:varmax-rule}
\end{equation}
where $g(z)$ is the varMax confidence statistic that combines the spread of the logits with their peak value, and $\tau$ is a calibrated threshold chosen on a held-out validation split that contains both in-distribution PAs and synthetic perturbations. This produces a five-way output: PA2, PA3, PA4, PA8, or \texttt{unknown}. We also retain the continuous confidence value $g(z)$ itself, which we expose to the downstream policy so that ``low-confidence PA2'' and ``high-confidence PA2'' induce distinct policy states.

For the application reported in~\cite{trott2026model}, we retain the network architecture validated on the pilot OTA dataset and adapt only the final decision layer. This isolates the effect of the open-set rule from architectural choices and lets us attribute downstream gains to the novelty-aware decision boundary rather than to representation changes.

\subsection{DQN-based Response Selection}
\label{subsec:dqn}

The varMax output is a label and a confidence, not an action. To map outputs to actions we adopt the DQN-IDS formulation of~\cite{tiwaridqn}, originally developed for open-set network intrusion detection, and adapt it to the EW response setting.

\paragraph{State.} At each decision step $t$, the agent observes
\begin{equation}
s_t = \bigl(\hat{y}_t,\ g(z_t),\ \phi_t,\ h_{t-1}\bigr),
\label{eq:state}
\end{equation}
where $\hat{y}_t \in \{\text{PA2, PA3, PA4, PA8, \texttt{unknown}}\}$ is the varMax verdict, $g(z_t)$ is its confidence, $\phi_t$ is the auxiliary event-statistics vector from preprocessing, and $h_{t-1}$ is a short history summary of recent verdicts (rolling counts and dwell time per class) so the policy can recognize behavioral patterns that span multiple windows. Note that DQN can functionally be an alternative to be used in the place of varMax. For such cases, we will compare the two algorithms isolated later in this paper.

\paragraph{Action space.} The agent selects an action $a_t \in \mathcal{A}$ from a discrete set of operationally meaningful responses. In our deployment $\mathcal{A}$ contains: continue passive observation; request a higher-rate capture on the same band; allocate the SDR's directional beam to localize the emitter; deploy an EMS deception waveform consistent with the framework of the parent project; and escalate to operator review. The set is small enough for tabular Q-learning to be infeasible only because of the continuous confidence component, which is why a function-approximation DQN is appropriate.

\paragraph{Reward.} The reward function is the place where QR-CWoS enters training. Rather than rewarding raw classification accuracy, we reward the agent for \emph{opening and maintaining} a usable window. For each episode the reward is the sum of (i)~a positive term proportional to the duration between PA detection and adversary tactic completion (the open window), (ii)~a negative term penalizing actions that unnecessarily consume sensing or deception resources, and (iii)~a large negative terminal reward if the adversary's tactic completes without the defender having selected a response other than passive observation. The \texttt{unknown} verdict is \emph{not} directly penalized; what is penalized is reaching the end of the episode having ignored an unknown.

\paragraph{Training.} We train the DQN with experience replay and a target network, following~\cite{tiwaridqn}. Episodes are constructed by drawing trajectories from the OTA dataset and interleaving held-out PAs at controlled rates, so the policy experiences both familiar and unseen tactics during training. To prevent the policy from overfitting to a particular adversary profile, we vary the inter-PA timing distribution and the prevalence of each PA across episodes.

\subsection{QR-CWoS Quantification}
\label{subsec:qrcwos}

The pipeline outputs a verdict and an action at every analysis window, but the operational quantity of interest is the QR-CWoS itself. We define the window for an episode as the interval $[t_d, t_c]$, where $t_d$ is the time at which the pipeline first emits a non-passive response with confidence above a deployment threshold, and $t_c$ is the time at which the adversary's full tactic would have completed in the absence of intervention (this completion time is recoverable in our OTA experiments because we control the emitter). Three derived quantities are reported:

\begin{enumerate}
    \item \textbf{Window duration}, $\Delta t = t_c - t_d$, the time during which the defender has acted but the adversary has not yet completed, i.e., the raw width of the QR-CWoS.
    \item \textbf{Window reliability}, the fraction of episodes for which $\Delta t$ exceeds a stated operational minimum, conditioned on whether the PA was seen or unseen at training time. This is the resilience axis of QR-CWoS: a wider mean window is not useful if the variance is high.
    \item \textbf{Window quality under uncertainty}, the conditional distribution of $\Delta t$ given the varMax confidence $g(z_t)$ at $t_d$. This is the quantification axis: it tells the operator how much window to expect for a given confidence level, and it lets the DQN's value function be interpreted as an estimate of expected window value.
\end{enumerate}

Reporting all three quantities rather than only mean accuracy is what distinguishes a QR-CWoS evaluation from a conventional classification evaluation. It is also what makes the pipeline auditable: an operator can ask, for any deployment, ``what window does this system actually deliver against tactics it has not seen, and how reliably?'', and the pipeline will answer in seconds rather than in classes.

\section{Experimental Design}
\label{sec:experiments}
We design the experimental protocol to test the routing problem introduced in Fig.~\ref{fig:dqnguard_pipeline}: given a closed-set RF preliminary-action backbone, can an OSR decision layer retain useful known-PA evidence while rejecting withheld PA behaviors as \texttt{unknown}? The experiments separate decision-rule effects from calibration-information effects by evaluating fixed-surrogate, target--surrogate, and target-open calibration regimes.

\subsection{Research Questions}
\label{subsec:rqs}

The experiments are organized around three questions:
\begin{itemize}
    \item \textbf{RQ1.} Under a fixed surrogate-open calibration condition, does DQNGuard improve open-set PA detection relative to VarMax and a DQN-IDS-style confidence head?
    \item \textbf{RQ2.} How sensitive is DQNGuard to the choice of surrogate-open calibration class?
    \item \textbf{RQ3.} When a small number of target-open examples are available for calibration, how much of the surrogate-open performance gap is explained by the absence of target-specific open evidence?
\end{itemize}

RQ1 evaluates the proposed decision rule at a common operating point. RQ2 tests whether surrogate-open calibration transfers across PA behaviors. RQ3 is diagnostic: it does not assume that target-open examples are available at deployment, but it helps determine whether failures are caused by poor surrogate choice rather than by the backbone representation alone.

\subsection{OTA Dataset and Capture Setup}
\label{subsec:dataset}

The OTA corpus contains WiFi, Bluetooth, and Zigbee emissions and five PA behaviors: Scan, Burst, Sustain, Hop, and Replay. These behavior names are used throughout the paper even though legacy repository identifiers are retained in code and result files for traceability.

Each spliced OTA window contains 400,000 complex I/Q samples. At the recorded OTA sample rate of 12.5 MS/s, each window spans 32 ms. For each protocol/action pair, the generation target is 10,000 windows. Burst, Sustain, Hop, and Replay were generated together per protocol and split across 20 shards. Scan was generated separately under the same seed schedule and capture setup and split across five shards. The shard structure is used for reproducibility and data management; the OSR evaluation is defined by the PA-class folds described below.

The OTA capture setup used two Ettus USRP N210 radios equipped with SBX/CBX daughterboards and VERT2450 omnidirectional antennas. The radios were separated by approximately 11 inches and connected through a NETGEAR ProSAFE GS108 Gigabit switch to a Lambda Vector Pro workstation. The captures used a 2.437 GHz center frequency with TX/RX gains fixed at 30/10 dB.

\subsection{Backbone Training and Splits}
\label{subsec:backbone_training}

Each fold designates some PA classes as known and withholds one class as the target unknown. The closed-set backbone is trained only on the known classes for that fold. The target unknown is excluded from closed-set training and appears only during open-set evaluation, except in the target-open diagnostic regime where a small labeled subset is used for OSR calibration.

The backbone produces the outputs consumed by all decision heads: logits, softmax probabilities, predicted known labels, and intermediate features. VarMax, DQN-IDS-style heads, and DQNGuard are therefore compared as decision layers over the same closed-set backbone outputs.

The evaluated runs use the OTA feature cache with cache length 16,384, source type \texttt{ota}, protocol tag \texttt{all}, dataset tag \texttt{ota\_core\_high\_run01}, learning rate $2\times10^{-4}$, entropy loss weight 0.05, and seed 0. Known samples are partitioned using the open-PA split mode into 70\% training, 15\% validation, and 15\% testing. Where calibration caps are used, the OSR calibration set contains 625 known samples and 625 open samples. The main DQNGuard operating point uses a 5\% known-rejection budget.

\subsection{Open-Set Calibration Regimes}
\label{subsec:calibration_regimes}

We evaluate three calibration regimes. The regimes differ only in what open evidence is available to the OSR decision layer; the closed-set backbone remains trained only on known PA classes.

\paragraph{Regime A: fixed-surrogate method comparison.}
The first regime fixes the surrogate-open calibration class to Scan and evaluates Burst, Sustain, Hop, and Replay as target unknowns. This setting isolates the OSR decision mechanism. Because the compared decision heads receive the same fixed surrogate-open information, performance differences are attributable primarily to the decision rule rather than to a different surrogate choice.

\paragraph{Regime B: Target--Surrogate Matrix.}
The second regime evaluates all ordered target/surrogate pairs. In each cell, one class is withheld as the target unknown, one different class is used as the surrogate-open calibration class, and the remaining three classes form the known set. With five PA classes, this produces 20 target/surrogate cells. This regime tests whether surrogate-open calibration is interchangeable across PA behaviors or whether surrogate usefulness is target-dependent.

\paragraph{Regime C: target-open diagnostic calibration.}
The third regime uses a small labeled subset of the held-out target-open class during OSR calibration, while still excluding that class from closed-set backbone training. This is not the assumed deployment condition. It is a diagnostic upper-bound condition that asks whether surrogate-open failures can be explained by the absence of calibration examples from the true target unknown.

\subsection{Metrics}
\label{subsec:metrics}

We report known rejection rate, unknown precision, unknown recall, unknown F1, OSR macro F1, OSR accuracy, and unknown AUROC. The known rejection rate is the fraction of known test samples incorrectly rejected as \texttt{unknown}. Unknown F1 measures binary detection of the held-out target-open class. OSR macro F1 averages performance across the known classes and the unknown class, penalizing methods that improve unknown detection by sacrificing known-class recognition.

Known rejection is central to the evaluation because unknown detection can be made artificially high by rejecting too many known samples. For DQNGuard, the main operating threshold is selected from known calibration scores only. With a 5\% known-rejection budget, the threshold is set at the 95th percentile of known calibration unknown scores. Open calibration samples may shape or diagnose guard evidence, but in the known-budgeted setting they do not determine the final operating threshold.


\section{Results}
\label{sec:results}
The results are organized around three questions: whether the proposed DQNGuard layer preserves known-class recognition while detecting unknown preliminary actions, how surrogate-open calibration transfers across target unknown classes, and how much performance changes under more permissive calibration assumptions.

\subsection{Main Open-Set Performance}

Table~\ref{tab:main_osr_results} compares the main OSR decision layers under the fixed PA1-surrogate known-budget calibration setting. DQNGuard achieves the highest unknown F1 and OSR macro F1 while maintaining a low known-rejection rate, indicating that the proposed decision layer improves unknown preliminary-action detection without obtaining the gain by simply rejecting known samples.

\begin{table}[t]
\centering
\caption{Main open-set recognition comparison under fixed PA1-surrogate known-budget calibration. Values are mean $\pm$ standard deviation across held-out preliminary-action folds.}
\label{tab:main_osr_results}
\scriptsize
\setlength{\tabcolsep}{3pt}
\begin{tabular}{lcccc}
\hline
Method & Known rej. $\downarrow$ & Unknown F1 $\uparrow$ & OSR macro F1 $\uparrow$ & AUROC $\uparrow$ \\
\hline
DQNGuard & 0.050 $\pm$ 0.004 & 0.865 $\pm$ 0.142 & 0.881 $\pm$ 0.109 & 0.891 $\pm$ 0.129 \\
Shreyash CNN + DQNGuard head & 0.063 $\pm$ 0.017 & 0.701 $\pm$ 0.197 & 0.764 $\pm$ 0.119 & 0.829 $\pm$ 0.092 \\
VarMax surrogate-all & 0.128 $\pm$ 0.043 & 0.745 $\pm$ 0.146 & 0.778 $\pm$ 0.104 & 0.951 $\pm$ 0.060 \\
\hline
\end{tabular}
\end{table}


\subsection{Target--Surrogate Transfer}

Figure~\ref{fig:target_surrogate_matrix} reports the Target--Surrogate Matrix for surrogate-open calibration. Each row corresponds to the true held-out unknown preliminary action, each column corresponds to the surrogate-open class used during calibration, and each cell reports unknown-class F1. The diagonal is omitted because the target unknown cannot also serve as its own surrogate.

\begin{figure}[t]
    \centering
    \includegraphics[width=\columnwidth]{figures/target_surrogate_matrix/target_surrogate_unknown_f1_matrix.pdf}
    \caption{Target--Surrogate Matrix for surrogate-open calibration. Cells report unknown-class F1, with bold entries marking the best surrogate for each target. Surrogate utility is highly target-dependent, showing that withheld classes are not interchangeable open-set calibration proxies.}
    \label{fig:target_surrogate_matrix}
\end{figure}

The matrix shows that surrogate choice materially changes unknown detection. For example, Scan is best detected when Hop is used as the surrogate-open class, whereas Sustain, Hop, and Replay are best detected when Scan is used as the surrogate. This asymmetry indicates that surrogate-open calibration does not provide a universal proxy for target-open behavior. Instead, surrogate usefulness depends on the relationship between the withheld calibration behavior and the true unknown behavior.


\section{Discussion}
\label{sec:discussion}
The results support three main observations.

First, open-set RF preliminary-action recognition should be evaluated as an operating-point problem, not only as a ranking problem. Table~\ref{tab:main_osr_results} shows that DQNGuard improves unknown F1 and OSR macro F1 while maintaining low known rejection under the fixed surrogate-open condition. This distinction matters because an OSR layer can appear effective if it rejects many samples, but such behavior would reduce the value of the known PA evidence stream. DQNGuard's known-budgeted thresholding makes the comparison stricter: unknown detection improves while the allowed loss of known PA utility remains constrained.

Second, surrogate-open calibration is useful but target-dependent. Figure~\ref{fig:target_surrogate_matrix} shows that the surrogate class is not an interchangeable source of generic ``unknownness'' evidence. Some surrogate behaviors induce guard evidence that transfers well to a target unknown, while others provide little useful separation. This suggests that surrogate-open calibration should be treated as a deployment assumption to be tested, not as a universal substitute for target-open calibration.

Third, the RF front end should be interpreted as a sensing and triage component. DQNGuard does not infer a complete attack chain, select an EW response, or assign final semantic labels to all novel emissions. Its narrower role is to decide whether an OTA RF window should be emitted as known PA evidence or preserved as a nonconforming behavior for later analysis. That boundary is important: the value of the method is not that it solves the downstream QR-CWoS policy problem, but that it supplies a cleaner evidence interface to such systems.

The main limitation is surrogate selection. The Target--Surrogate Matrix demonstrates that some surrogates transfer better than others, but the current experiments do not provide a reliable automatic rule for choosing the best surrogate before the target unknown is observed. A practical deployment would therefore require either a conservative fixed-surrogate policy, pooled surrogate evidence, human-reviewed surrogate diagnostics, or a learned surrogate-selection rule trained across a broader behavior taxonomy.

A second limitation is scale. The evaluated corpus spans three protocol families and five PA behaviors, which is sufficient to expose target-dependent surrogate transfer but not sufficient to characterize the full space of RF behavior novelty. Larger PA taxonomies would allow leave-multiple-out evaluation, pooled calibration studies, and stronger tests of whether the observed transfer asymmetries are specific to the current behaviors or reflect a broader property of RF open-set recognition.

Finally, the evaluation begins after RF window extraction and feature construction. Operational performance will also depend on upstream burst detection, segmentation, synchronization, and OTA capture quality, as well as downstream label-making and response-policy modules. These dependencies define the boundary of the claim. The result shown here is that unknown PA behaviors can be detected from OTA RF windows under controlled calibration regimes, while preserving a bounded stream of known PA evidence.


\section{Conclusion}
\label{sec:conclusion}
This paper introduced DQNGuard, a budgeted open-set decision layer for RF preliminary-action recognition from OTA captures. DQNGuard uses predicted-class calibration, variance/energy-style guard evidence, and known-only thresholding to decide whether each RF window should be retained as known PA evidence or routed as an unknown behavior.

Across WiFi, Bluetooth, and Zigbee OTA captures, DQNGuard improves the fixed-budget operating point relative to VarMax and a DQN-IDS-style confidence head. The Target--Surrogate Matrix further shows that surrogate-open calibration is target-dependent: withheld surrogate classes are not interchangeable proxies for future unknown behaviors.

These results support RF open-set recognition as a sensing and triage layer for QR-CWoS rather than as a complete response system. Future work will focus on principled surrogate selection, pooled calibration, larger RF behavior taxonomies, and downstream attack-chain integration.


\bibliographystyle{IEEEtran}
\bibliography{references}


\end{document}
